\documentclass[a4paper,12pt]{article}
\usepackage[utf8]{inputenc}
\usepackage[T1]{fontenc}
\usepackage[margin=2.5cm]{geometry}
\usepackage{ctex} % 中文支持
\usepackage{graphicx}
\usepackage{booktabs} % 三线表支持
\usepackage{xcolor}
\usepackage{enumitem} % 列表定制
\usepackage{float}
\usepackage{titlesec}

% -----------------------------------------------------------
% 格式设置区：Level 1 高结构化 + Level 2 标准专业风格
% -----------------------------------------------------------

% 设置层级标题格式（自动编号，移除手动序号）
% 使用稳重的深蓝色/深灰色，符合专业报告调性
\titleformat{\section}{\Large\bfseries\color{blue!40!black}}{\thesection}{1em}{}
\titleformat{\subsection}{\large\bfseries\color{gray!80!black}}{\thesubsection}{1em}{}

% 设置段落间距与缩进
\setlength{\parskip}{0.6em}
\setlength{\parindent}{0pt}

% 设置列表间距
\setlist{nosep} % 紧凑列表

% -----------------------------------------------------------
% 文档元数据
% -----------------------------------------------------------
\title{\textbf{2024年度环境、社会及管治（ESG）报告：可持续发展管理与绩效}}
\author{本公司董事会办公室}
\date{2025年3月}

\begin{document}

\maketitle

% 前言
\section*{前言}

2024年，本公司坚持将环境、社会及管治（ESG）理念融入企业发展战略与日常运营管理。作为国内领先的化工行业上市公司，我们严格遵守《环境保护法》、《安全生产法》及相关监管要求，致力于通过技术创新与精细化管理，实现经济效益与社会效益的协同发展。本章节将重点披露公司在关键ESG议题上的管理实践与年度绩效。

% 第一章 环境
\section{环境管理：合规运营与低碳转型}

公司建立了完善的环境管理体系，坚持源头减量、过程控制与末端治理相结合的原则，切实履行环境保护责任。

\subsection{环境关键绩效指标 (KPIs)}

报告期内，公司加强了环境数据的监测与统计。2024年度主要环境绩效数据如下：

\begin{table}[H]
\centering
\caption{2024年度环境排放与资源利用绩效表}
\begin{tabular}{@{}ll@{}}
\toprule
\textbf{指标类别} & \textbf{2024年度数据} \\ \midrule
\textbf{排放管理} & \\
温室气体排放总量 & 39,107.65 吨CO2e \\
排放强度 & 4.58 吨/百万元营收 \\
一般工业固废产生量 & 170.48 吨 \\
危险废物合规处置率 & 100\% \\
\midrule
\textbf{资源利用} & \\
新鲜水总用量 & 149.55 万吨 \\
包装材料使用总量 & 7,677 吨 \\
\midrule
\textbf{环保投入} & \\
环保资金投入 & 457.93 万元 \\
缴纳环保税 & 9.6 万元 \\ \bottomrule
\end{tabular}
\end{table}

\subsection{清洁生产与节能降耗行动}

为响应国家“双碳”战略，公司在生产及运营环节实施了多项技术改造，重点推进能源结构的优化与循环经济的实践。

\begin{itemize}
    \item \textbf{热能清洁替代（核心生产基地）}：报告期内，该基地完成了供热系统的清洁化改造，通过引入清洁能源热源替代原有的燃煤锅炉设施。该项目的实施有效消除了燃烧环节的废气排放，实现了生产过程废气的“零排放”，符合国家关于工业锅炉整治的合规要求。
    \item \textbf{水资源循环利用（化工厂）}：公司旗下化工厂实施了冷凝水回收利用项目。通过回收生产蒸汽冷凝水作为锅炉补水，实现了热能与水资源的双重回收，提升了资源利用效率。
\end{itemize}

\subsection{绿色矿山与运输电气化}

针对移动源排放管理，公司在\textbf{西北某大型露天煤矿项目}积极推进运输设备的电气化升级，以降低化石能源消耗。

\begin{enumerate}
    \item \textbf{设备更新}：本年度新增21台纯电动矿卡及36台油电增程式矿卡。新型纯电动矿卡具备制动能量回收功能，显著降低了运营能耗。
    \item \textbf{智能化应用}：投入12台无人驾驶电动矿卡，通过自动驾驶技术优化行驶路径，实现了作业过程的零排放与低噪音，减少了对矿区周边生态环境的影响。
\end{enumerate}

% 第二章 创新
\section{研发创新：驱动行业高质量发展}

技术创新是公司核心竞争力的体现，也是推动化工与民爆行业安全、高效发展的关键动力。

\subsection{研发资源与成果}

公司持续保持高强度的研发投入，构建了完善的知识产权保护体系。

\begin{itemize}
    \item \textbf{研发投入}：2024年研发支出共计4.19亿元，占营业收入比例为4.91\%。
    \item \textbf{人才队伍}：截至报告期末，公司拥有研发人员1,064人。
    \item \textbf{知识产权}：本年度新增授权专利126项，累计拥有有效专利596项。
\end{itemize}

\subsection{核心技术与应用价值}

公司重点推广以下三项具有行业领先水平的技术，致力于提升本质安全水平：

\begin{enumerate}
    \item \textbf{现场混装水胶炸药工艺}：
    \begin{itemize}
        \item \textbf{技术特点}：该技术实现了原材料运输与混合工序的分离，填补了国内技术空白，达到国际先进水平。
        \item \textbf{应用价值}：通过现场即时混合，消除了成品炸药运输过程中的安全风险；同时具备密度可调特性，能够根据地质条件优化爆破效果。
    \end{itemize}
    
    \item \textbf{二氧化碳相变膨胀激发管}：
    \begin{itemize}
        \item \textbf{技术特点}：国内首创技术，产能规模达500万只。利用液态二氧化碳物理相变原理破岩。
        \item \textbf{应用价值}：全过程无明火产生，显著提升了高瓦斯矿井等复杂环境下的开采安全性。
    \end{itemize}
    
    \item \textbf{智能矿山作业系统}：
    \begin{itemize}
        \item \textbf{技术特点}：集成5G通讯、边缘计算与机器人技术。
        \item \textbf{应用价值}：实现了5G钻机、自动验孔机器人与无人矿卡的协同作业，有效降低了高危岗位的人员暴露风险，推动矿山作业向无人化、智能化转型。
    \end{itemize}
\end{enumerate}

% 第三章 社会
\section{社会责任：员工权益与社区共建}

公司坚持“以人为本”的管理理念，切实保障员工权益，并积极投身公益事业，回馈社会。

\subsection{职业健康与安全管理 (OHS)}

公司始终将安全生产置于首位，构建了全方位的安全保障体系。截至2024年，员工劳动合同签订率保持100\%。

\begin{table}[H]
\centering
\caption{2024年度安全保障投入及成效}
\begin{tabular}{@{}ll@{}}
\toprule
\textbf{项目} & \textbf{说明/金额} \\ \midrule
安全生产责任险投入 & 355.99 万元 \\
工伤保险投入 & 892.5 万元 \\
应急管理体系 & “1+14+N”应急预案体系 \\ \bottomrule
\end{tabular}
\end{table}

\textbf{海外项目安全实践}：
在\textbf{纳米比亚}矿服项目中，公司克服跨文化管理挑战，建立了标准化的安全管理制度（如“晨会制度”及“手指口述”确认法）。报告期内，该项目实现\textbf{连续130万小时无损工事件}，展现了公司国际化运营中的安全管理水平。

\subsection{企业公民责任}

公司积极响应国家乡村振兴战略，并在突发灾害救援中发挥专业优势。

\begin{itemize}
    \item \textbf{公益投入}：2024年对外捐赠及公益支出总额1,219万元，其中乡村振兴专项投入1,125万元。
    \item \textbf{应急救援案例}：在2024年\textbf{湖南某地}决堤抢险行动中，公司迅速响应，调集专业技术力量。通过采用“松动爆破”技术方案，在确保石料完整性的前提下，短时间内开采石料8.4万吨，为抗洪抢险提供了关键的物资保障。
\end{itemize}

% 第四章 治理
\section{公司治理：合规与商业道德}

作为\textbf{上级集团}成员企业，公司致力于完善治理结构，维护投资者权益，并打造廉洁透明的供应链。

\subsection{治理结构与信披质量}

\begin{itemize}
    \item \textbf{董事会架构}：公司构建了“3+3+3”的董事会结构（3名股东董事、3名独立董事、3名执行董事），确保决策的科学性与制衡性。
    \item \textbf{监管评价}：凭借规范的运作与透明的披露，公司在证券交易所年度信息披露考评中获得“A级”评价。
\end{itemize}

\subsection{股东回报与供应链管理}

\begin{itemize}
    \item \textbf{利润分配}：公司重视股东投资回报，2024年度派发现金红利2.54亿元（每10股2.05元），分红比例占可分配利润的40.12\%。
    \item \textbf{廉洁从业}：公司对供应链实施严格的商业道德管理。针对1,548家合作供应商，推行《廉洁协议》签署制度，签署率达到100\%，有效防范商业贿赂风险。
\end{itemize}

\end{document}